% ==================== 摘要专用页 ====================
\begin{center}
  {\heiti\zihao{2} 微网与外部电网电力调控策略}
\end{center}

\vspace{0.6em}

\noindent{\heiti\zihao{4} 摘\quad 要}

\vspace{0.3em}

本文针对微网在光伏本地消纳与外网购电成本之间的权衡问题，围绕\textbf{储能设备的优化调度}，
建立了一套由确定性到随机、由单时点到多时点的递进式数学模型，并给出四问的完整求解结果。
全文统一采用 10 分钟为时段（一天 144 段），储能容量 $[1200,10800]$ kWh、
单向效率 $0.9$、功率上限 $5000$ kW，不售电且允许弃光，缺口按 5 倍电价紧急购电。

\textbf{问题一}为典型日确定性调度。以全天购电费最小为目标，建立含 720 个决策变量的
稀疏线性规划（LP），用 HiGHS 求解。我们证明了\textbf{该 LP 的松弛是紧的}：当往返效率
$0.9^2=0.81<1$ 时，若某时段同时充电与放电，则可构造等 SOC 扰动使购电量严格下降，
故最优解自动满足充放电互斥，\textbf{无需引入二元变量}。求解得全天最优购电费
\textbf{35\,126.95 元}，较无储能基线 48\,052.05 元\textbf{节省 26.90\%}；其中光伏余电消纳贡献
6\,727.28 元、峰谷价差套利贡献 6\,197.82 元，两者几乎各占一半。

\textbf{问题二}为全年滚动随机日前调度。我们首先识别出本题的\textbf{报童型结构}：计划购电量照付
（缺货成本为 0、过量成本为电价 $p$），而缺额须按 $5p$ 补救，故临界分位
$F^*=c_u/(c_o+c_u)=0.80$，即计划量应取净负载的 \textbf{80\% 分位}而非均值。据此建立
\textbf{两阶段随机规划}：用「计划净供给 $\ge$ 点预测」的不等式约束叠加历史残差场景的紧急购电
兜底，配合 48 小时前瞻窗口逐日滚动。预测严格只用历史数据（负载：最近 4 次同星期均值
$+$ 近 5 日残差漂移；光伏：近 7 日均值 $+$ 近 28 日残差漂移）。关键改进是
\textbf{剔除 1 月 1--7 日的回退预测残差}（该 7 天负载内核缺少同星期历史，口径与后续不一致），
以 56 天等权残差构造场景，全年总费用 \textbf{1\,438.91 万元}。

\textbf{问题三}引入 0/6/12/18 点的光伏发布预报与多时点调整机制。首先将题面「计划量高于调整量
部分按 50\% 违约电价、调整量高于计划量部分按 1.5 倍」的计费规则统一为
$J=\sum_t[p_tg_t^0+1.5p_tu_t-0.5p_tv_t+5p_te_t]$，并指出\textbf{调减项的系数必须为 $-0.5p$}
（净省 $0.5p$）——否则最优解永不调整，题目失去意义。其次提出\textbf{在线预报融合}：
以历史预测 $H$ 与发布预报 $F$ 的因果最小二乘权重 $\alpha$ 融合，并做收缩偏差校正。
进而构建\textbf{经验场景树}：把历史日在后续发布时刻的预报修订作为未来信息做二叉聚类，
\textbf{节点内共享购电/充放电/SOC 变量}，从而在变量结构上严格保证非前视。
全年费用 \textbf{1\,402.24 万元}，较同口径基线\textbf{节省 2.55\%}，紧急购电量减少 28.74\%。
边际分析表明 6 点与 12 点调整各贡献约 8 万元，\textbf{18 点几乎无额外收益}。
进一步通过完美预见下界（1\,223.16 万元）与权重饱和性检验等 30 余组对照，
说明该方案在给定数据与已测试模型族内已接近性能上界。

\textbf{问题四}考虑实时波动电价。对 4-2，我们保留「当天 0 点电价未知」的主模型，
以 1 月历史选定「最近 4 次同星期几均值」为电价预测器，并构造
\textbf{电价—净负载联合场景}（同一历史日的两种残差配对），
理由是 $\E[p_t(N_t-z_t)_+]\neq\E[p_t]\E[(N_t-z_t)_+]$——高电价可能与缺电同时发生，
不能分解。全年费用 \textbf{1\,516.30 万元}。对 4-3，在 4-2 上叠加多时点调整与预报融合机制，
费用进一步降至 \textbf{1\,506.17 万元}。对比可见\textbf{波动电价下多时点调整的价值
（约 10.1 万元）显著大于固定电价下的价值（约 1.5 万元）}，说明电价波动放大了调整机制的价值。

\vspace{0.4em}
\noindent{\heiti 关键词：}\quad 储能优化调度；线性规划；报童模型；两阶段随机规划；
滚动时域优化；预报融合；场景树；联合场景

\vspace{0.8em}
\noindent\rule{\textwidth}{0.4pt}
